
\section{Sandeep}

A critical question in understanding the labor market and informing policy recommendations is whether the observed differences in earnings in the formal and informal employment is the result of market segmentation or the incentives within competitive market framework.  In other words, the question is whether individuals opt for informal employment by choice or are pushed into it due to the barriers of entry into formal employment. Although the question is particularly important in developing economies, where the phenomenon of informality is widespread; no consensus has been achieved on the driving mechanism of informality due to the complex structure of informal employment.

The traditional dual market theory, starting with (Lewis, 1954) conceive informal employment as backward, marginal and disadvantaged sector where employees are paid wages above market-clearing prices and serves as a last resort to escape unemployment (Harris and Todaro, 1970, FIelds, 1990, 2005a, Stighlitz, 1976). According to this approach, workers in the informal sector is predominantly involuntarily engagement  and earn less than identical workers in the formal sector. This segmentation may result from access to information or geographic barriers to entry, or from sector-specific variations in employment and wage regulations (Dikens and lang, 1992). For instance, firms may offer wages above market clearing prices to increase profit by increasing retention, reducing turnover costs, eventually maximizing productivity of employees, raise price of labor and reduce the demand for workers in the formal sector.      

The notion of informal employment as a transitional ground for entry into formal employment was challenged by (ILO, 1973), which saw informal sector as permanent, economically sufficient and profit-making small-scale urban economic continuum of employment and income, with little evidence of holding onto sectoral mobility (Sethuraman, 1981). This evidence gave rise to an alternative theory which states that informal employment results from the voluntary choices of workers within the competitive market framework (Maloney, 1999, Rosenzweig, 1988). Working in the informal sector reflects worker's own choice, based on the cost-benefit analysis of respective sector, given their individual characteristics. Voluntary employment in the informal sector can be due to the workers trading their pecuniary rewards for non-pecuniary characteristics where individuals maximise their utility rather than their earnings (Maloney, 2004),  when there is possibility of self-selection into jobs based on individual's comparative advantage (Roy, 1951 ) or motivated by the desire to avoid tax burdens and inefficient regulations (Fields, 2003).  

Hence, two theories explain informal employment: the segmentation hypothesis views it as a survival strategy for those excluded from formal jobs, while the comparative advantage hypothesis sees it as a voluntary choice by workers to maximize their utility or earnings. While the differential between formal and informal employment is well-documented (Mazumdar, 1975, Pradhan and Van Soest 1995, Badaoui et al. 2008, Blunch, 2015, ), lower wages and returns to education and experience in informal employment doesnot necessarily imply market segmentation as long as there is mobility between these two sectors (Gunther and Launov, 2012).  The observed differences in wage distributions is associated with nonrandom selection process or served as a penalty due to differences in worker's skill or abilities. It is therefore necessary to account for the unobservable characteristics of workers which determines sector choices and earnings in order to distinguish between segmentation and comparative advantage. A range of empirical studies have explored whether labor markets in developing and developed countries operate competitively or are segmented (Dickens and lang, 1985, Heckman and Hotz, 1986, Rosenzweig, 1988, Magna, 1991, Pratap and Quintin, 2006). %Have to review this paragraph based on the need whether it fits our results well . 

The empricial literature on informal employment is predominantly based on the assumption that the informal sector is homogeneous (Gunther and Launov, 2012). This assumptions has been With the seminal works of Fields (1990), theoretical studies a heterogeneous informal sector with two distinct tiers:  \textit{lower-tier} and \textit{upper-tier}. The  \textit{lower-tier} fits the dual maket model assuming a formal employment paying high wages and providing job security and social security contributions, whereas and an informal employment assuming individuals who cannot afford unemployment, paid low wages and are excluded from the formal jobs. The \textit{upper-tier} comprises individuals who voluntarily chose informal employment, expected to earn more than formal employment due to their specific characteristics and non-job benefits(Cunningham and Maloney, 2001, Henley et. al., 2009). 













In Figure 1, we present estimates of the corresponding quantities from the quantile decompositions across the overall distribution of wage gap attributable to differences in the average characteristics of the two groups (composition effect) and differences in returns to those characteristics (structure effect). These estimates are detailed for the overall service sector as well as its stratified segments—upper tier, middle tier, and lower tier across two distinct years, 2008 and 2018. The stratified segment reveal heterogeneity in the relative importances of covariates in accounting for disparities across the distribution, which the unidimensional classical method cannot capture. The results reveal that the wage gap is maximum below the median ($\tau$ 0.25 - 0.5), whereas it deteriorates at an increasing rate above.  





%Figure \vref{fig:service} shows the decomposition result of the formal-informal occupation-based wage gap for the service sector. We observed an improvement in the wage gap from 2008 to 2018. However, the gap is still prominent for the lower wage quantile. In both periods, the gap increased until it nears the median point, and we observed a sharp decline while moving beyond the median point towards a higher wage quantile. This nature of the wage gap is explained more by the structure effect (SE), which also depicts a similar nature to a total effect, than the composition effect (CE) where . We observed a slight decline in CE, only for those below the median, but overall, it is constant throughout the distribution. \par 


\section{Sanjeet}
The level of wages and earnings of the workforce is determined by the forces of demand for and supply of labor based on the neo-classical framework. The same however is not true for developing countries, where labor market dualism and strong entry barriers across different segments of the labor market exists. The dual labor market model supposes the existence of two distinct sectors of economic activity usually classified as formal and informal sector, which extends to the nature of the employment \footnote{The distinction between the formal and informal sector on the one hand and formal and informal employment on the other has become increasingly important with growing numbers of employed persons who work in formal sector establishments also do not have access to basic benefits. The focus on informal employment is therefore to identify persons who work in precarious employment situations, irrespective of the sector into which the entity for which they work falls.} . Within formal employment, employees are subject to stable jobs with higher pay, better working conditions, permanent contract, paid parental and annual leave, paid leave for illness and accidents and social security contributions provided by the employer. Although informal employment has no employment contracts and does not ensures these jobs based benefits, it does provide more flexibility in terms of the working hour and places of work. This asymmetry affects how employees value formal and informal employment differently. 

Although those in poverty are more likely to participate in informal employment, their wage rates and total income from informal work are lower than formal works, which furthermore reinforces their marginalized position (Williams, 2014). The productivity level of informal employment is lower than the formal employment due to the differences in technology \textcite(Monsted2020) (Monsted, 2020), lack of transparency of their own accounts, long established territory based on well-entrenched territory control and rents, sub-optimal allocation of productive factors, reliance on family sources for credit, prevents acquiring modern management skills and worker training, limiting growth potential and access to world market and fragile working condition (Benjamin et al., 2014). Because of the mismatch in productivity and subsequently different wage setting mechanism, the extent of wage discrimination for the formal and informal employees is likely to vary based on personal and household, job and employer characteristics. Within formal employment, there is anti-discrimination and equal treatment laws and regulations in place, whereas these regulations and wage bargaining process are completely absent for informal employment, with competition as well as social norms favoring the dominant group (Neog \& Sahoo, 2016). Understanding of the wage discrimination with the existence of distinct employment opportunities and differing characteristics via distinct policy recognition is crucial to understand the labor market dynamics and welfare distribution within the economy.

This paper focuses on the individuals of 15-60 years who are in full time wage employment defined as individuals as who work in private financial business firm, private non-financial institutions, non-profit institutions and other institutions for 40 hours or more during the week with cash as the mode of the payment \footnote{Those who work in government, state-owned enterprises and international organizations/foreign embassy are excluded. Similarly, formal sector and informal sector both comprises of non-agriculture data and those whose economic activities are in the agricultural industries are reported on separately within job sectors.} . It aims to examine whether significant gaps in earnings exists among both formal and informal wage earners. There are several reasons motivating for the choice of full time wage employees. First, it is important to analyze employees separately from other groups, such as self-employed workers, because certain determinants of differences in earnings such as employer characteristics and job characteristics apply only to full time wage employees. Second, the definition of informality is not universal and the definition varies across employment statuses. In my study, I have used the conceptual framework of informal employment as defined by the 15th International Conference of Labour Statisticians as the Nepal Labor Force Survey III follows the same definition \footnote{15th International Conference of Labour Statisticians (ILO, 2018) defines Informal Employment as “the total number of informal jobs, whether carried out in formal sector enterprises, informal sector enterprises, or households, or as the total number of persons engaged in informal jobs during a given reference period”. In the Nepal Labour Force Survey III, those in informal employment are identified and the residual is those in formal employment.} . Third, focusing on employees makes it easier to compare and identify the gaps in literature which is concentrated mostly on gender wage gap and lacks the study based on wage gaps on informality. Fourth, excluding agriculture sector, 41 percent of total employment is based on informal sector, which compromises the single largest most sector of employment in Nepal. To complement the analysis. I provide the analysis of formal and informal employment based on 

In this paper, I use OLS regressions to assess the roles of individual and household, firm and employer characteristics in shaping hourly wages of formal and informal employees. To recover the gaps in potential wage offers, I control for these characteristics of formal and informal employment. Moreover, I use Oaxaca-Blinder Decomposition to determine whether employees in informal employment faces a wage gap in comparison to formal employment. A participation decision of workers to participate in any sector (endogenous selection) brings selectivity bias to occur and affects the formal–informal wage gap. Failing to deal with non-random selection into formal and informal jobs is a major concern as it would lead to misleading estimates of wage gaps. In order to solve the problem of sample selection, I introduce a correction term obtained from a probit model into the wage equation. Using the available data, I expect the proportion of the informal sector to all workers to influence the sectoral choice of workers without having a direct impact on wages. The Oaxaca-Blinder methodology breaks down variations in average outcomes between groups into differences in average characteristics and rewards to those characteristics. Furthermore, I carry out quantile regression to capture workers heterogeneity in returns to observed characteristics which allows us to estimate wage gap at different quantiles of wage distribution. Finally, I use Machado and Mata (MM) decomposition with a bootstrap approach, extending the traditional Oaxaca-Blinder wage decomposition on the mean across the entire wage distribution combining quantile regression and bootstrap approach to explain the significance of covariates across all distribution. Since the unconditional quantile is not the same as the integral of the conditional quantiles, authors provide a simulation-based estimator where the counterfactual distribution is constructed from the generation of a random sample.

Using the data from the nationally representative labor force surveys conducted in 2017 referred as Nepal Labor Force Survey- third round (NLFS -III), I find that those in formal employment have greater average wages than those in informal employment. 

The remainder of this paper is organized as follows: Section 2 presents the development of our econometric model. In Section 3, we describe the dataset and provide descriptive statistics on wage inequality among formal and informal employees in Nepal. Section 4 discusses the results, focusing on selection into OLS-RIF regression and decomposition based on quantile regression for formal and informal employment in Nepal. Finally, Section 5 concludes the paper with potential avenues for further research.   




\section{Chapter Plan}

1. Introduction 
2. Econometric Model (Cherkonov paper) 
- 
3. Empirical application
4. Conclusion
5. Appendix
6. Biblography

1. Introduction
2. Data
3. Empirical Application
4. Results and Discussion
5. Conclusion


Occupation/ Industry
HH formal
Occupation proportion


Econometric model
- The original method of decomposing differences in the mean of an outcome variable has been extended to distributional parameters other than the mean (Juhn, Murphy and Pierce (1993) and DiNardo, Fortin and Lemieux(1996).   

The heterogeneity of informal employment in developing country shows that it arises from the individual skills of workers in respective sectors, comparatitive advantage from working formally vs informally but also the mechanism inducing workers to work informally (Radchenko, 2016). This raises a question on the suitability of dual market vs integrated market classification, implying that rationing frontiers between two theories are hard to explain. 


The Oaxaca-Blinder decomposition serves as a foundational method for understanding mean wage differentials between two groups (i.e formal and informal employment), further divided into explained and unexplained effect in the literature. This is consisent with the empirical approach of estimating Mincerian wage equations by least squares methods, which provide estimates of the effect of covariates upon the mean of the conditional wage distribution \parencite{Machado2005}. A significant body of methodological work has been added to the literature aimed at refining and expanding the analysis to include distributional parameters beyond the mean.  \textcite{Machado2005} relies on quantile regressions for each possible quqntile combined with the simulation procedure, \textcite{DiNardo1996} suggest a reweighting procedure to compute distribution of dummy covariates to the aggregate composition effect. In contrast to these path dependent methods, \textcite{Firpo2009} uses the recentered influence function for the distribution statistic of interest in a regression. 


The decomposition of the overall differences into two components depends on the construction of meaningful counterfactual wage distribution based on the partial equillibrium approah trying to answer the counterfactual question "What wage would a person employed in the informal employment have if s/he was instead employed in a similar job in the formal employment?".  %Counterfactuals used in the decompositions often consist of manipulating structural wage setting function linking the observed and unobserved characteristics of workers to their wages for each group. Since this simple counterfactuals may not be appropriate for answering the economic question of interest, there has been a development of differing methods for estimation of counterfactuals. For instance, the formal employment wage structure may not represent the appropriate counterfactual for the way employees are paid in informal employment in the absence of labor market discrimination.
Several alternative methods to construct counterfactual marginal quantile functions were reviewed in \textcite{Fortin2011}. \textcite{DiNardo1996} uses a non-parametric reweighting approach using the propensity score, \textcite{Machado2005} and (Autor, Katz and Kearney, 2005) use conditional quantile regression based estimation procedure, whereas \textcite{Firpo2009} uses a recentered influence function regression approach. In our study, we use  \textcite{Chernozhukov2013}'s distribution regression procedure to derive a couterfactual marginal cumulative distribution which is inverted to provide quantile gap estimates for estimating the semiparametric estimators of conditional distribution.  

Suppose we would like to analyze the wage differences between formal and informal employment. Let $i$ and $f$ denote the population of informal and formal employment respectively. The variable $Y_j$ denotes wages and $X_j$ denotes job-market relevant characteristics that affect the wages for populations $j = i$ and $j = f$. The conditional distribution functions $F_{Y_i|X_i} (y|x)$ and $F_{Y_f|X_f} (y|x)$ describe the stochastic assignment of wages to workers with characteristics $x$, for informal and formal employment respectively.  Let  $F_{Y(i|i)}$ and $F_{Y(f|f)}$ represent the observed distribution of wages for informal and formal employment respectively, and let $F_{Y({f|i})}$ represent the distribution function of wages that have prevailed for informal employment had they faced the formal employment's wage schedule $F_{Y_f|X_f}$ , known as the counterfactual distribution constructed by integrating the conditional distribution of wages for formal employment with respect to the distribution of characteristics for informal employment : 

\begin{linenomath*}
	\begin{align}
		F_{Y(i|f)} (y)= \int_{X_1} F_{Y_i|X_i}(y|x) \, . dF_{X_f}(x) 
	\end{align}
\end{linenomath*} 

With the understanding that selection could be different for different wage equations in the market, we propose an estimator for subgroup of individuals that doesnot require any prior information about selection based on (Machado, 2018). We divide the group as "always employed", which are individuals who work in formal or informal employment regardeless of instruments.

In the literature, four primary approaches address selection concerns in wage gap analyses \parencite{Machado2017}. The first approach utilizes information on observed covariates to impute missing data, assuming that after controlling for detailed characteristics, selection into employment is random. Methods include matching similar individuals and using panel data to fill in missing wages \parencite{Blau2006, Neal2004, Olivetti2016}. The second approach posits that selection bias diminishes when participation rates are high, allowing for estimates from smaller samples with nearly complete participation, although this limits generalizability \parencite{Mulligan2008} , (Chamberlain, 1986). The third approach models the self-selection process by incorporating a control function in the wage equation to account for unobservable factors influencing selection, which often necessitates an exclusion restriction \parencite{Heckman1974}, (Vella, 1998). However, this method's sensitivity to modeling assumptions can affect reliability \parencite{Bar2015} (Melly, 2015). The fourth approach employs economic theory to impose bounds on the worst-case scenario for the wage gap, reducing reliance on unobserved data through instruments that shift participation \parencite{Manski1990, Blundell2007}. While this method imposes fewer restrictions, it often requires additional assumptions to produce informative bounds.






